\section{Implementation}
In this section, we want to show how the design choices made in the design section (\ref{sec:sys-design}) have been implemented in the code. 
\section{Read}
For the read operation, the implementation was straightforward. The replica simply returns the value associated with the requested key from its positions array.
Thanks to the broadcast protocol's design, a \texttt{ReadRequest} message sent after reciving a \texttt{WriteResult} for a previous
write operation will always return the value of the write or a more recent one.
The \texttt{ReadRequest} message are handled by the \texttt{onReadRequest} method of the \texttt{Replica} class.
Even during the election protocol, the \texttt{ReadRequest} messages are handled in the same way.
While during an election, the replica may not have the most up-to-date information, it will still return a possibly outdated but still valid
\texttt{ReadResult} message, avoiding uncessary read timeouts for the client.
\section{Write}
% TODO @AleCornella
\section{Election Acknowledgment}
% TODO @ozneroL541
